
\subsection{The earlier reconstruction applied to every original source level}
\label{sec:support-propagated-original-levels}

We now evaluate (D2), (D6), and (D7) on the original theta source,
including the finite polynomial source levels used in (TC2) and
(TB.24)--(TB.26).  Let $\mathcal W$ contain every complex linear
subspace $W$ of the original $V$, ordered by inclusion; no closure
or invariance under $D$ is imposed for this cochain construction.
Fix $n\geq1$ and $I_n=\{0,\ldots,n-1\}$.  Its nonbottom labels are
\[
\lambda=(W,S;A),\qquad W\in\mathcal W,
\quad\varnothing\ne S\subseteq\{+,-\},\quad A\subseteq I_n.
\tag{D9}
\]
Adjoin one distinct bottom $\bot$, with zero fibre.  The join of two
nonbottom labels is $(W+W',S\cup S';A\cup A')$.  The bottom is
strictly below even $(0,S;\varnothing)$.  In degrees zero and one put
\[
\begin{split}
C_{W,S;A}^{X,0}&=
 \left(\bigoplus_{+\in S}W\ \oplus\!
       \bigoplus_{-\in S}\widehat W\right)^A,
\qquad C_{W,S;A}^{X,1}=\mathscr B_X^A,\\
d_{W,S;A}(\phi_i,\psi_i)_i&=(\Theta\phi_i-\Theta\widehat\psi_i)_i.
\end{split}
\tag{D10}
\]
Here $X=U$ or $\Omega$ as in (MCF41); missing legs use typed zeros.
The transition includes old subspaces and coordinates and inserts
typed zeros in newly present legs or coefficients.  Fourier
linearity gives $\widehat W\subseteq\widehat{W+W'}$; hence all
transitions have their stated domains.  Substitution into the
differential proves the cochain identity, and inclusions compose
literally.  Thus (D2) applies in each degree to this actual diagram.
At every nonbottom label with $A=\varnothing$ the amplitude fibre
is zero, but its element is $(W,S;\varnothing,0)$, while the global
zero is $(\bot,0)$.  Applying (D2) proves they remain distinct and
gives their exact joins.  No coefficient-mask colimit is taken.

For each $W$, put $Q_{W,X}=\mathscr B_X/\Theta W$.  The full
cohomology and the comparison to the full source are
\[
\begin{gathered}
H^0(C_{W,S;A}^X)=
 \begin{cases}W^A,&S=\{+,-\},\\0,&|S|=1,\end{cases}
\qquad H^1(C_{W,S;A}^X)=Q_{W,X}^A,\\
0\longrightarrow(V/W)^A\xrightarrow{\eta_W^A}Q_{W,X}^A
 \xrightarrow{\pi_W^A}(T_XQ)^A\longrightarrow0,
\qquad \eta_W[v]=[\Theta v]_W.
\tag{D11}
\end{gathered}
\]
Indeed the source differential has image $(\Theta W)^A$.  A
single-leg differential is injective.  At a joint slot the inverse
coordinates $u=\phi-\widehat\psi$, $v=(\phi+\widehat\psi)/2$,
$\phi=v+u/2$, $\psi=\widehat v-\widehat u/2$ identify its kernel
with $v\in W$.  For the last sequence, injectivity of $\Theta$
makes $\eta_W$ injective; a class killed by $\pi_W$ is represented
by $\Theta v$ for some $v\in V$, and every element of $T_XQ$ has
a representative in $\mathscr B_X$.  This proves exactness,
coordinatewise also for empty $A$.
The coequalizer (D6) uses precisely $(\Theta W)^A$, and (D7)
identifies the kernel of the full-source comparison with $(V/W)^A$
through this same $\eta_W^A$.  This is the all-label fibre kernel
of (D8); with unchanged support label, the global-zero categorical
kernel contains only its bottom fibre.  Both maps are retained.

For $W\subseteq W'$ the source transition therefore induces
\[
0\longrightarrow(W'/W)^A\xrightarrow{[v]\mapsto[\Theta v]_W}
 Q_{W,X}^A\longrightarrow Q_{W',X}^A\longrightarrow0.
\tag{D12}
\]
The proof uses the same kernel computation with $W'$ in place of
$V$.  Thus source refinement has its actual nonzero comparison
kernel before any full-source quotient is applied.

The critical-strip enlargement also has its original typed map:
\[
0\longrightarrow C_{W,S;A}^{U}\longrightarrow C_{W,S;A}^{\Omega}
 \longrightarrow(T_LQ)^A[-1]\longrightarrow0,
\qquad L=\{\Re s=1/2\}.
\tag{D13}
\]
The first map is identity on source legs and inclusion on the
central term.  Its degree-zero quotient is zero.  Its degree-one
quotient is $\mathscr B_\Omega^A/\mathscr B_U^A$, since the same
$\Theta W$ lies in both.  The quotient map $q$ identifies this
space with $(T_\Omega Q/T_UQ)^A$.  For any annihilator
$p=p_Up_L$ with its factors rooted in the disjoint sets $U,L$,
choose $ap_U+bp_L=1$.  The maps $b(D)p_L(D)$ and $a(D)p_U(D)$
split each killed class into its $U$ and $L$ parts.  Their
intersection is zero by the same Bezout identity; uniqueness makes
the splitting independent of the annihilator.  Thus the quotient
is exactly $(T_LQ)^A$, with every multiplicity retained.  All
three complexes keep the same outer and coefficient labels.

The original arithmetic operator has the following additional
type on this larger source diagram.  If $DW\subseteq W'$, then
\[
D_{W,W'}:C_{W,S;A}^X\longrightarrow C_{W',S;A}^X,
\quad(\phi_i,\psi_i)_i\longmapsto(D\phi_i,(1-D)\psi_i)_i,
\quad(F_i)_i\longmapsto(DF_i)_i.
\tag{D14}
\]
Fourier gives $(1-D)\widehat W=\widehat{DW}\subseteq\widehat W'$,
and $\Theta D=D\Theta$ proves the cochain identity.  Moreover
$D\mathscr B_X\subseteq\mathscr B_X$ because $D$ commutes with
each polynomial annihilator of a quotient class.  Thus every
displayed map has its original domain.  For $D$-invariant $W$,
take $W'=W$ to recover precisely (MCF34).  For the actual finite
source levels $W_m=\operatorname{span}\{\phi_*,D\phi_*,\ldots,D^m\phi_*\}$,
one instead has $DW_m\subseteq W_{m+1}$ and must use (D14) with
that codomain.  Linear independence of these vectors follows from
$H_{P(D)\phi_*}=P$, so $\dim(W_{m+1}/W_m)=1$ exactly.
Dilation likewise gives a cochain map from $W$ to $\mathcal R_aW$,
using $\mathcal R_a$ on the plus source and central term and
$a\mathcal R_{1/a}$ on the minus source.  The Fourier identity
proves its stated minus codomain $\widehat{\mathcal R_aW}$.
